\chapter{Homology}

\begin{center}\itshape
Count the vertices of a cube, subtract its edges, add its faces: $8-12+6=2$. Do the same for a tetrahedron, an icosahedron or a soccer ball and you get
$2$ again. Drill a tunnel through the solid and you get $0$. Homology explains the miracle: the alternating count does not see the faces at all, only the
holes, counted dimension by dimension. With it we prove that an open set of $\R^n$ cannot be an open set of $\R^m$ unless $m=n$, and that at every moment
two opposite points of the Earth have the same temperature and the same pressure.
\end{center}

On 14 November 1750, Leonhard Euler wrote to Christian Goldbach from Berlin about a discovery he found astonishing for its simplicity: in every solid
bounded by plane faces, the number of solid angles plus the number of faces exceeds the number of edges by two. He could not prove it in general, and
the proof he published in 1758, by slicing off corners one at a time, has a gap. Legendre gave a proof by spherical geometry in 1794, Cauchy another by
flattening the polyhedron into a plane network in 1813.

Then the counterexamples arrived. In 1813 Simon Lhuilier, a mathematician in Geneva, pointed out solids for which the formula fails: a picture frame,
whose central tunnel makes $V-E+F=0$; a cube with a smaller cube sitting on top of one face, whose top face is a square with a hole; a solid with a
hollow inside. Each time the definition of ``polyhedron'' was adjusted to exclude the monster, and each time a new monster appeared. In the 1960s the
philosopher Imre Lakatos turned this history into \emph{Proofs and Refutations}, a dialogue in a classroom where students rediscover every
counterexample, and argued that mathematics grows exactly this way: proofs suggest the right definitions, and counterexamples suggest the right proofs.

The right definition was homology. Poincar\'e's Betti numbers counted independent cycles: closed curves that do not bound a piece of surface, closed
surfaces that do not bound a solid, and so on. For thirty years they were numbers, manipulated with matrices. In the mid-1920s Emmy Noether, in
G\"ottingen, insisted in lectures and conversations that the cycles modulo the boundaries form a \emph{group}, and that the Betti numbers and the torsion
coefficients are just its invariants. Heinz Hopf used her viewpoint in 1928 to give a short algebraic proof of the Euler--Poincar\'e formula; Pavel
Alexandroff and Hopf recorded her influence in the preface of their \emph{Topologie} of 1935. Every theorem of this chapter is a statement about
these groups.

We build singular homology, which is defined for every space and is obviously a topological invariant, and simplicial homology, which is defined by
finite matrices and is obviously computable, and we prove that they agree. Then we compute spheres and surfaces, prove invariance of domain and the
Borsuk--Ulam theorem, and identify the Euler characteristic.

\section{Chain complexes}

\begin{definition}
A \emph{chain complex} $C$ is a sequence of abelian groups and homomorphisms $\cdots\to C_{k+1}\xrightarrow{\partial}C_k\xrightarrow{\partial}C_{k-1}\to\cdots$ with $\partial\circ\partial=0$, and $C_k=0$ for $k<0$.
Elements of $Z_k=\ker(\partial\colon C_k\to C_{k-1})$ are \emph{cycles}, elements of $B_k=\partial(C_{k+1})$ are \emph{boundaries}, and the \emph{homology} is $H_k(C)=Z_k/B_k$. A \emph{chain map}
$f\colon C\to C'$ is a family of homomorphisms with $\partial f=f\partial$; it maps cycles to cycles and boundaries to boundaries, hence induces $f_*\colon H_k(C)\to H_k(C')$. Two chain maps
$f,g$ are \emph{chain homotopic} if there are $P\colon C_k\to C'_{k+1}$ with $\partial P+P\partial=g-f$; then $f_*=g_*$, because for a cycle $z$, $g(z)-f(z)=\partial P(z)$ is a boundary.
\end{definition}

A sequence of homomorphisms is \emph{exact} if the image of each map is the kernel of the next. A \emph{short exact sequence} of chain complexes $0\to A\xrightarrow{i}B\xrightarrow{j}C\to0$ consists of chain
maps with $0\to A_k\to B_k\to C_k\to0$ exact for each $k$.

\begin{lemma}[the zigzag lemma]\label{lem:zigzag}
A short exact sequence of chain complexes $0\to A\xrightarrow iB\xrightarrow jC\to0$ gives a long exact sequence
\[\cdots\to H_k(A)\xrightarrow{i_*}H_k(B)\xrightarrow{j_*}H_k(C)\xrightarrow{\delta}H_{k-1}(A)\xrightarrow{i_*}H_{k-1}(B)\to\cdots\to H_0(C)\to0,\]
where $\delta[c]=[a]$ for $c=j(b)$ and $i(a)=\partial b$. It is natural: a commutative diagram of short exact sequences gives a commutative diagram of long exact sequences.
\end{lemma}

\begin{proof}
\emph{$\delta$ is defined.} Given a cycle $c\in C_k$, write $c=j(b)$; then $j(\partial b)=\partial c=0$, so $\partial b=i(a)$ for a unique $a$, and $i(\partial a)=\partial\partial b=0$ gives $\partial a=0$. Another choice $b'=b+i(a')$
changes $a$ by $\partial a'$, and replacing $c$ by $c+\partial c'$ with $c'=j(b'')$ changes $b$ by $\partial b''$ (plus an element of $i(A)$), which does not change $\partial b$. So $\delta[c]=[a]$ is well defined, and clearly a homomorphism.

\emph{Exactness at $H_k(B)$.} $j_*i_*=0$ since $ji=0$. If $j(b)=\partial c'$ with $c'=j(b')$, then $j(b-\partial b')=0$, so $b-\partial b'=i(a)$, and $a$ is a cycle because $i$ is injective; so $[b]=i_*[a]$.

\emph{At $H_k(C)$.} $\delta j_*[b]=0$ since $\partial b=0$. If $\delta[c]=0$, then $a=\partial a'$, and $b-i(a')$ is a cycle with $j(b-i(a'))=c$.

\emph{At $H_{k-1}(A)$.} $i_*\delta[c]=[\partial b]=0$. If $i(a)=\partial b$, then $c=j(b)$ is a cycle ($\partial c=j\partial b=ji(a)=0$) with $\delta[c]=[a]$.

Naturality: each map in the definition of $\delta$ commutes with the given chain maps.
\end{proof}

\begin{lemma}[the five lemma]\label{lem:five}
In a commutative diagram of abelian groups with exact rows
\[\begin{array}{ccccccccc}A_1&\to&A_2&\to&A_3&\to&A_4&\to&A_5\\ \downarrow\alpha_1&&\downarrow\alpha_2&&\downarrow\alpha_3&&\downarrow\alpha_4&&\downarrow\alpha_5\\ B_1&\to&B_2&\to&B_3&\to&B_4&\to&B_5\end{array}\]
if $\alpha_1,\alpha_2,\alpha_4,\alpha_5$ are isomorphisms, so is $\alpha_3$.
\end{lemma}

\begin{proof}
Write $f_i$, $g_i$ for the horizontal maps out of $A_i$, $B_i$. \emph{Injective:} if $\alpha_3(x)=0$, then $\alpha_4f_3(x)=0$, so $f_3(x)=0$ and $x=f_2(y)$. Then $g_2\alpha_2(y)=0$, so
$\alpha_2(y)=g_1(z')=g_1\alpha_1(z)=\alpha_2f_1(z)$, hence $y=f_1(z)$ and $x=f_2f_1(z)=0$. \emph{Surjective:} given $w\in B_3$, $g_3(w)=\alpha_4(u)$; then $\alpha_5f_4(u)=g_4g_3(w)=0$, so $f_4(u)=0$ and $u=f_3(x)$.
Now $g_3(w-\alpha_3(x))=0$, so $w-\alpha_3(x)=g_2(\alpha_2(y))=\alpha_3(f_2(y))$, and $w=\alpha_3(x+f_2(y))$.
\end{proof}

\section{Simplicial homology}

\begin{definition}
A (finite, ordered) \emph{simplicial complex} $K$ is a finite totally ordered set of vertices and a family of nonempty subsets, the \emph{simplices}, containing all singletons and all
nonempty subsets of its members. A simplex with $k+1$ vertices has dimension $k$ and is written $[v_0,\dots,v_k]$ with $v_0<\dots<v_k$. Its \emph{realization} is the convex hull of the
corresponding basis vectors of $\R^{\text{vertices}}$, and $|K|$ is the union of these, a compact space. A \emph{subcomplex} is a subfamily that is itself a simplicial complex. A
\emph{triangulation} of a space $X$ is a homeomorphism $|K|\to X$.
\end{definition}

Let $R$ be $\Z$, $\Q$ or $\Z/2$. The group $C^\Delta_k(K;R)$ of simplicial $k$-chains is the free $R$-module on the $k$-simplices, and
\[\partial[v_0,\dots,v_k]=\sum_{i=0}^k(-1)^i[v_0,\dots,\widehat{v_i},\dots,v_k],\]
the hat meaning omission. Each $(k-2)$-face $[\dots\widehat{v_i}\dots\widehat{v_j}\dots]$, $i<j$, appears twice in $\partial\partial$, with signs $(-1)^j(-1)^i$ (removing $v_j$ first) and $(-1)^i(-1)^{j-1}$ (removing
$v_i$ first, which shifts the position of $v_j$), so $\partial\partial=0$. The \emph{simplicial homology} is $H^\Delta_k(K;R)=H_k(C^\Delta(K;R))$; for a subcomplex $L$, the relative homology $H^\Delta_k(K,L)$ is the
homology of the quotient complex $C^\Delta(K)/C^\Delta(L)$.

\begin{figure}[ht]
\centering
\fig{ch06-torus}
\caption{A triangulation of the torus: a $3\times3$ grid of squares, each cut by a diagonal, with opposite sides of the big square glued (the vertex labels show the
gluing). It has $V=9$, $E=27$, $F=18$, and $V-E+F=0$. A coarser grid would not do: with a $2\times2$ grid, two different triangles would have the same three vertices.}
\label{fig:torustriang}
\end{figure}

For coefficients $\Z$, the boundary maps are integer matrices, and homology is computed by the Smith normal form of Volume~2, Theorem~7.4.

\begin{proposition}\label{prop:smithhomology}
Let $c_k$ be the number of $k$-simplices, $r_k$ the rank of $\partial\colon C_k\to C_{k-1}$, and $d_1\mid\dots\mid d_{r_{k+1}}$ the nonzero invariant factors of $\partial\colon C_{k+1}\to C_k$. Then
\[H^\Delta_k(K;\Z)\cong\Z^{b_k}\oplus\bigoplus_{d_i>1}\Z/d_i,\qquad b_k=c_k-r_k-r_{k+1}.\]
Consequently (the Euler--Poincar\'e formula)
\[\chi(K):=\sum_k(-1)^kc_k=\sum_k(-1)^kb_k.\]
\end{proposition}

\begin{proof}
By the Smith normal form, there are bases $(f_j)$ of $C_{k+1}$ and $(e_i)$ of $C_k$ with $\partial f_i=d_ie_i$ for $i\le r=r_{k+1}$ and $\partial f_j=0$ for $j>r$. Then $B_k=\bigoplus_{i\le r}d_i\Z e_i$. For $i\le r$,
$d_i\partial e_i=\partial\partial f_i=0$, so $\partial e_i=0$: the $e_i$, $i\le r$, are cycles, and $Z_k=\bigoplus_{i\le r}\Z e_i\oplus Z'$ with $Z'=Z_k\cap\operatorname{span}(e_{r+1},\dots)$, a free group whose rank is
$\operatorname{rank}Z_k-r=(c_k-r_k)-r_{k+1}$. So $H_k=Z_k/B_k\cong\bigoplus_{i\le r}\Z/d_i\oplus Z'$. In the alternating sum of the $b_k$, the ranks $r_k$ cancel in pairs.
\end{proof}

The $b_k$ are the \emph{Betti numbers}. With coefficients in a field $R=\Q$ or $\Z/2$ the same proof, with Gaussian elimination instead of the Smith form, gives $H^\Delta_k(K;R)\cong R^{c_k-r_k-r_{k+1}}$ (ranks
computed over $R$), and the same Euler characteristic. The torus of Figure~\ref{fig:torustriang} has $\chi=9-27+18=0$; the boundary of a tetrahedron has $\chi=4-6+4=2$. Nothing yet says that
two triangulations of the same space have the same $\chi$; that will follow from singular homology.

\section{Singular homology}

Let $\Delta^n$ be the standard simplex, the convex hull of the basis vectors $e_0,\dots,e_n$ of $\R^{n+1}$. A \emph{singular $n$-simplex} in a space $X$ is a map $\sigma\colon\Delta^n\to X$. The group $C_n(X;R)$ is the
free $R$-module on all singular $n$-simplices, an enormous group, and
\[\partial\sigma=\sum_{i=0}^n(-1)^i\,\sigma|[e_0,\dots,\widehat{e_i},\dots,e_n],\]
where $\sigma|[w_0,\dots,w_{n-1}]$ denotes $\sigma$ composed with the affine map $\Delta^{n-1}\to\Delta^n$ sending $e_j$ to $w_j$. The computation that gave $\partial\partial=0$ for simplicial chains applies word for
word. The \emph{singular homology} $H_n(X;R)$ is the homology of this complex; we drop $R$ when it is $\Z$. A map $f\colon X\to Y$ induces $f_\#\sigma=f\circ\sigma$, a chain map, so $f_*$ is defined on homology,
with $(gf)_*=g_*f_*$: homeomorphic spaces have isomorphic homology.

For a subspace $A\subset X$, $C(A)$ is a subcomplex of $C(X)$, and the \emph{relative homology} $H_n(X,A)$ is the homology of $C(X)/C(A)$; a relative cycle is a chain whose boundary lies in $A$. By
the zigzag lemma there is a long exact sequence of the pair,
\[\cdots\to H_n(A)\to H_n(X)\to H_n(X,A)\xrightarrow{\delta}H_{n-1}(A)\to\cdots,\]
natural for maps of pairs. The \emph{reduced homology} $\tilde H_n(X)$ is the homology of the complex augmented by $\varepsilon\colon C_0(X)\to R$, $\varepsilon(\sigma)=1$; it equals $H_n(X)$ for $n>0$, and
$H_0(X)\cong\tilde H_0(X)\oplus R$ for $X\ne\emptyset$.

\begin{proposition}\label{prop:H0}
$H_0(X;R)$ is free on the path components of $X$. If $X$ is a point, $\tilde H_n(X)=0$ for all $n$.
\end{proposition}

\begin{proof}
A $0$-simplex is a point, and a $1$-simplex is a path $\sigma$ with $\partial\sigma=\sigma(e_1)-\sigma(e_0)$. So two points are homologous exactly when they are joined by a chain of paths, and $H_0$ is free on
the path components. For a point, there is one singular simplex $\sigma_n$ in each dimension, and $\partial\sigma_n=\sum_i(-1)^i\sigma_{n-1}$ is $\sigma_{n-1}$ for $n$ even and positive, $0$ for $n$ odd. The augmented
complex $\cdots\to R\xrightarrow{1}R\xrightarrow0R\xrightarrow1R\xrightarrow0R\xrightarrow{\varepsilon}R$ is exact.
\end{proof}

\begin{theorem}[homotopy invariance]\label{thm:homotopyH}
If $f,g\colon X\to Y$ are homotopic, then $f_*=g_*$. Consequently a homotopy equivalence induces isomorphisms in homology, and a contractible space has $\tilde H_*=0$.
\end{theorem}

\begin{proof}
Let $F\colon X\times I\to Y$ be the homotopy. In $\Delta^n\times I$ let $v_i=(e_i,0)$ and $w_i=(e_i,1)$. The \emph{prism operator} is
\[P(\sigma)=\sum_{i=0}^n(-1)^i\,F\circ(\sigma\times\mathrm{id})|[v_0,\dots,v_i,w_i,\dots,w_n]\in C_{n+1}(Y),\]
where $|[\dots]$ is again composition with the affine map from $\Delta^{n+1}$. Write $\Phi=F\circ(\sigma\times\mathrm{id})$. Then
\[\partial P\sigma=\sum_{j\le i}(-1)^{i+j}\Phi|[v_0,..,\widehat{v_j},..,v_i,w_i,..,w_n]+\sum_{j\ge i}(-1)^{i+j+1}\Phi|[v_0,..,v_i,w_i,..,\widehat{w_j},..,w_n].\]
In the terms with $j=i$, $\Phi|[v_0,..,v_{i-1},w_i,..,w_n]$ appears with sign $+$ from the first sum (index $i$) and with sign $-$ from the second sum (index $i-1$), so these cancel, except
$\Phi|[w_0,\dots,w_n]=g\circ\sigma$ from $i=0$ and $-\Phi|[v_0,\dots,v_n]=-f\circ\sigma$ from $i=n$. The terms with $j\ne i$ are exactly $-P(\partial\sigma)$, as one sees by expanding $P$ on each face of $\sigma$ (the face
missing $e_j$ contributes the prisms with $i<j$ with sign $(-1)^i(-1)^j$, and those with $i>j$ with the index shifted). So $\partial P+P\partial=g_\#-f_\#$, a chain homotopy.

If $gf\simeq\mathrm{id}$ and $fg\simeq\mathrm{id}$ then $g_*f_*=\mathrm{id}$ and $f_*g_*=\mathrm{id}$. A contractible space is homotopy equivalent to a point.
\end{proof}

The prism operator maps chains in a subspace $A$ to chains in any subspace containing $F(A\times I)$; so homotopy invariance holds for maps of pairs $f,g\colon(X,A)\to(Y,B)$ homotopic through maps of
pairs. With the five lemma, if $A\subset X$ is a deformation retract then $H_n(X,A)=0$, and if $A\subset A'$ is a deformation retract of $A'$ then $H_n(X,A)\to H_n(X,A')$ is an isomorphism.

\section{Small chains and Mayer--Vietoris}

The key technical fact is that homology can be computed with arbitrarily small simplices. Let $\mathcal U=\{U_j\}$ be a family of subsets of $X$ whose interiors cover $X$, and $C^{\mathcal U}(X)$ the
subcomplex of chains made of singular simplices each lying in some $U_j$.

\begin{theorem}[small chains]\label{thm:smallchains}
The inclusion $\iota\colon C^{\mathcal U}(X)\to C(X)$ induces isomorphisms in homology.
\end{theorem}

\begin{proof}
\emph{Linear chains.} In a convex set $Y\subset\R^N$, call \emph{linear} the singular simplices $[w_0,\dots,w_n]$ that are affine maps, and $LC(Y)$ the chains they span. For a point $b$, the \emph{cone}
$b[w_0,\dots,w_n]=[b,w_0,\dots,w_n]$ satisfies $\partial(b\lambda)=\lambda-b(\partial\lambda)$ for $n\ge1$, and $\partial(b\lambda)=\lambda-\varepsilon(\lambda)[b]$ for $n=0$. Define the \emph{barycentric subdivision}
$S\colon LC_n(Y)\to LC_n(Y)$ by $S=\mathrm{id}$ on $0$-chains and $S\lambda=b_\lambda(S\partial\lambda)$ for a simplex $\lambda$ with barycentre $b_\lambda$, and $T\colon LC_n\to LC_{n+1}$ by $T=0$ on $0$-chains and
$T\lambda=b_\lambda(\lambda-T\partial\lambda)$. By induction on $n$, $\partial S=S\partial$ and $\partial T+T\partial=\mathrm{id}-S$: for instance $\partial T\lambda=\lambda-T\partial\lambda-b_\lambda(\partial\lambda-\partial T\partial\lambda)$, and by induction
$\partial T\partial\lambda=\partial\lambda-S\partial\lambda-T\partial\partial\lambda=\partial\lambda-S\partial\lambda$, so $\partial T\lambda=\lambda-T\partial\lambda-S\lambda$. By induction, every simplex of $S\lambda$ has vertices $b_{F_0},\dots,b_{F_n}$, the
barycentres of a chain of faces $F_n\subset\dots\subset F_0=\lambda$. If $F\subset G$ are faces and $G$ has vertices $w_0,\dots,w_k$, then $|b_G-b_F|\le\max_i|b_G-w_i|$ (a convex function on the convex hull is
largest at a vertex) and $|b_G-w_i|=\frac1{k+1}|\sum_j(w_j-w_i)|\le\frac k{k+1}\operatorname{diam}G$. So each simplex of $S\lambda$ has diameter at most $\frac n{n+1}\operatorname{diam}\lambda$, and those of $S^m\lambda$
at most $(\frac n{n+1})^m\operatorname{diam}\lambda$.

\emph{General chains.} For $\sigma\colon\Delta^n\to X$, let $S\sigma=\sigma_\#S(\mathrm{id}_{\Delta^n})$ and $T\sigma=\sigma_\#T(\mathrm{id}_{\Delta^n})$, where $\mathrm{id}_{\Delta^n}$ is a linear simplex in the convex set $\Delta^n$. Since $S$
and $T$ on linear chains commute with affine maps (which preserve barycentres), these agree with the previous definitions on linear chains, and the identities $\partial S=S\partial$ and
$\partial T+T\partial=\mathrm{id}-S$ carry over (apply them to $\mathrm{id}_{\Delta^n}$ and push forward by $\sigma_\#$). Let $D_m=\sum_{0\le i<m}TS^i$; then $\partial D_m+D_m\partial=\mathrm{id}-S^m$. Each simplex of $T\sigma$ and of $S\sigma$ has
image inside the image of $\sigma$.

\emph{Choosing $m$.} For each $\sigma$, the sets $\sigma^{-1}(\operatorname{int}U_j)$ cover the compact $\Delta^n$; by the Lebesgue number and the diameter bound there is a least $m(\sigma)$ such that every simplex of
$S^{m(\sigma)}\sigma$ lies in some $U_j$; $m(\sigma)=0$ if $\sigma$ itself does, and $m(\tau)\le m(\sigma)$ for every face $\tau$ of $\sigma$. Define $D\sigma=D_{m(\sigma)}\sigma$ and $\rho=\mathrm{id}-\partial D-D\partial$. Then
\[\rho\sigma=S^{m(\sigma)}\sigma+\sum_\tau\pm\bigl(D_{m(\sigma)}-D_{m(\tau)}\bigr)\tau,\]
summed over the faces $\tau$ of $\sigma$, and
\[\bigl(D_{m(\sigma)}-D_{m(\tau)}\bigr)\tau=\sum_{m(\tau)\le i<m(\sigma)}TS^i\tau\]
is a chain in $C^{\mathcal U}(X)$, because the simplices of $S^i\tau$ lie in sets $U_j$ for $i\ge m(\tau)$ and $T$
keeps a simplex inside the image of the simplex it comes from. So $\rho$ maps $C(X)$ into $C^{\mathcal U}(X)$; it is a chain map, because $\partial\rho=\partial-\partial D\partial=\rho\partial$; and $\rho\iota=\mathrm{id}$ because $D=0$ on
$C^{\mathcal U}(X)$. Finally $\mathrm{id}-\iota\rho=\partial D+D\partial$. So $\iota_*$ and $\rho_*$ are inverse isomorphisms.
\end{proof}

\begin{theorem}[Mayer--Vietoris]\label{thm:mv}
If $X=\operatorname{int}A\cup\operatorname{int}B$, there is a long exact sequence
\[\cdots\to H_n(A\cap B)\xrightarrow{(i_{A*},\,i_{B*})}H_n(A)\oplus H_n(B)\xrightarrow{j_{A*}-j_{B*}}H_n(X)\xrightarrow{\delta}H_{n-1}(A\cap B)\to\cdots\]
with all maps induced by inclusions except $\delta$, and the same in reduced homology when $A\cap B\ne\emptyset$. If $z=x+y$ is a cycle with $x$ in $A$ and $y$ in $B$, then $\delta[z]=[\partial x]$.
\end{theorem}

\begin{proof}
The sequence $0\to C(A\cap B)\to C(A)\oplus C(B)\to C^{\{A,B\}}(X)\to0$, with maps $c\mapsto(c,c)$ and $(a,b)\mapsto a-b$, is exact: a chain in both $C(A)$ and $C(B)$ is a chain in $A\cap B$. Apply the
zigzag lemma and replace $H(C^{\{A,B\}}(X))$ by $H(X)$, using Theorem~\ref{thm:smallchains}. For reduced homology, augment each complex in degree $-1$ by $0\to R\to R\oplus R\to R\to0$.
\end{proof}

\begin{corollary}[excision]\label{cor:excision}
If $Z\subset A\subset X$ with the closure of $Z$ inside the interior of $A$, then the inclusion $(X\setminus Z,A\setminus Z)\to(X,A)$ induces isomorphisms $H_n(X\setminus Z,A\setminus Z)\cong H_n(X,A)$.
\end{corollary}

\begin{proof}
Let $B=X\setminus Z$; then $\operatorname{int}A\cup\operatorname{int}B=X$. The quotient $C^{\{A,B\}}(X)/C(A)$ equals $C(B)/C(A\cap B)$, since a chain made of simplices in $A$ or in $B$ is, modulo $C(A)$, a chain in $B$.
The inclusion $C^{\{A,B\}}(X)\to C(X)$ induces isomorphisms in homology (Theorem~\ref{thm:smallchains}) and is the identity on $C(A)$, so by the long exact sequences of $C(A)\subset C^{\{A,B\}}(X)$ and
$C(A)\subset C(X)$ and the five lemma, it induces isomorphisms $H(C^{\{A,B\}}(X)/C(A))\cong H_n(X,A)$.
\end{proof}

\section{Spheres, degree and the agreement of the two theories}

\begin{theorem}\label{thm:spheres}
$\tilde H_k(S^n;R)\cong R$ for $k=n$ and $0$ otherwise. More precisely, let $U_\pm$ be the open sets $\pm x_{n+1}>-\frac12$ of $S^n$ and $S^{n-1}$ the equator $x_{n+1}=0$. The Mayer--Vietoris map
$\delta\colon\tilde H_k(S^n)\to\tilde H_{k-1}(U_+\cap U_-)\cong\tilde H_{k-1}(S^{n-1})$ is an isomorphism, natural for maps $S^n\to S^n$ that preserve $U_+$, $U_-$ and the equator.
\end{theorem}

\begin{proof}
$U_\pm$ are contractible (they deformation retract onto the poles), and $U_+\cap U_-$ deformation retracts onto the equator $S^{n-1}$. In reduced Mayer--Vietoris the terms $\tilde H(U_\pm)$ vanish, so
$\delta\colon\tilde H_k(S^n)\cong\tilde H_{k-1}(S^{n-1})$ for all $k$. For $S^0=\{\pm1\}$, $\tilde H_0$ is $R$, generated by $[1]-[-1]$, and all other groups vanish. Induction finishes.
\end{proof}

A map $f\colon S^n\to S^n$, $n\ge1$, acts on $H_n(S^n;\Z)=\Z$ as multiplication by an integer, its \emph{degree}. By Theorem~\ref{thm:homotopyH} the degree is a homotopy invariant; $\deg(fg)=\deg f\deg g$;
$\deg\mathrm{id}=1$; a map that is not onto has degree $0$, since it factors through $S^n\setminus\{p\}\cong\R^n$, which is contractible. For $n=1$, the degree is the winding number, as we will see in
Example~\ref{ex:surfacesH}.

\begin{proposition}\label{prop:reflection}
A reflection of $S^n$, such as
\[r(x_1,\dots,x_{n+1})=(-x_1,x_2,\dots,x_{n+1}),\]
has degree $-1$, and the antipodal map has degree $(-1)^{n+1}$.
\end{proposition}

\begin{proof}
$r$ preserves $U_\pm$ and commutes with the Mayer--Vietoris sequence, so by Theorem~\ref{thm:spheres} its degree equals that of the same reflection of $S^{n-1}$, and so on down to $S^0$, where $r$
exchanges $1$ and $-1$ and so maps the generator $[1]-[-1]$ of $\tilde H_0$ to its negative. Reflections in other coordinate hyperplanes are conjugate to $r$ by a permutation of coordinates $h$, and
$\deg(hrh^{-1})=\deg r$. The antipodal map is the composite of the $n+1$ coordinate reflections.
\end{proof}

\begin{corollary}\label{cor:brouwerH}
There is no retraction of the closed ball $\bar B^{n+1}$ onto $S^n$, and every map of $\bar B^{n+1}$ to itself has a fixed point. If $m\ne n$ ($m,n\ge1$), no nonempty open subset of $\R^m$ is homeomorphic to an open
subset of $\R^n$.
\end{corollary}

\begin{proof}
A retraction $r$ with inclusion $i$ would give $\mathrm{id}=r_*i_*$ on $\tilde H_n(S^n)=\Z$, factoring through $\tilde H_n(\bar B^{n+1})=0$. The fixed point theorem follows as in Chapter~3
(Theorem~\ref{thm:brouwern}), now with no smoothness needed. If $U\subset\R^m$ and $V\subset\R^n$ are homeomorphic by $h$, and $x\in U$, then by excision (cut away the complement of a small ball)
$H_k(U,U\setminus\{x\})\cong H_k(\R^m,\R^m\setminus\{x\})$, which by the long exact sequence is $\tilde H_{k-1}(\R^m\setminus\{x\})=\tilde H_{k-1}(S^{m-1})$: nonzero exactly for $k=m$. The homeomorphism carries
$(U,U\setminus\{x\})$ to $(V,V\setminus\{h(x)\})$, so $m=n$.
\end{proof}

We now compare the two theories. For a finite ordered simplicial complex $K$, send each simplex $[v_0,\dots,v_k]$ to the affine singular simplex $\Delta^k\to|K|$ with $e_i\mapsto v_i$. This is a chain map
$\varphi\colon C^\Delta(K;R)\to C(|K|;R)$, compatible with subcomplexes.

\begin{theorem}\label{thm:simplicialsingular}
$\varphi_*\colon H^\Delta_k(K;R)\to H_k(|K|;R)$ is an isomorphism for every $k$, and likewise for pairs $(K,L)$.
\end{theorem}

\begin{proof}
Induction on the number of simplices. Suppose the theorem holds for complexes with fewer simplices than $K$, and let $\sigma$ be a simplex of $K$ of largest dimension $k$, and $L$ the subcomplex obtained by
removing it; all proper faces of $\sigma$ lie in $L$. Let $\bar\sigma$ be the complex of all faces of $\sigma$ and $\partial\bar\sigma$ that of its proper faces. By the five lemma applied to the long exact sequences
of the pairs $(K,L)$ (the zigzag lemma is natural with respect to $\varphi$), it is enough to know that $\varphi_*$ is an isomorphism for $L$ (true by induction) and for the pair $(K,L)$.

\emph{The simplex.} $H^\Delta(\bar\sigma)$ is $R$ in degree $0$ and $0$ elsewhere: if $v$ is the first vertex, $h(\tau)=[v,\tau]$ for $v\notin\tau$ and $h(\tau)=0$ for $v\in\tau$ satisfies $\partial h+h\partial=\mathrm{id}$ in positive degrees
(check on $\tau$ with and without $v$). Also $|\bar\sigma|$ is convex, so contractible. So $\varphi_*$ is an isomorphism for $\bar\sigma$, and by induction for $\partial\bar\sigma$ (fewer simplices); by the five lemma, for the
pair $(\bar\sigma,\partial\bar\sigma)$.

\emph{Comparing the pairs.} On the simplicial side, $C^\Delta(\bar\sigma)/C^\Delta(\partial\bar\sigma)$ and $C^\Delta(K)/C^\Delta(L)$ are both $R$, generated by $\sigma$ in degree $k$, so the inclusion is an isomorphism.
On the singular side, let $b$ be the barycentre of $\sigma$ and $U=|K|\setminus\{b\}$. Since $\sigma$ is not a face of another simplex, its interior is open in $|K|$, and $U$ and $|\sigma|$ have interiors covering $|K|$.
Radial projection from $b$ (the identity on $|L|$) shows that $|L|$ is a deformation retract of $U$ and $|\partial\sigma|$ one of $|\sigma|\setminus\{b\}$. So $H(|K|,|L|)\cong H(|K|,U)$ and
$H(|\sigma|,|\partial\sigma|)\cong H(|\sigma|,|\sigma|\setminus\{b\})$. Excision (Corollary~\ref{cor:excision}) used twice, removing $Z=|L|$ from $A=U$ in $|K|$ and $Z=|\partial\sigma|$ from $A=|\sigma|\setminus\{b\}$ in $|\sigma|$,
identifies both $H(|K|,U)$ and $H(|\sigma|,|\sigma|\setminus\{b\})$ with $H(\operatorname{int}\sigma,\operatorname{int}\sigma\setminus\{b\})$, all maps being induced by inclusions. So $(|\sigma|,|\partial\sigma|)\to(|K|,|L|)$ induces isomorphisms, and in the commutative
square formed with $\varphi_*$ three maps are isomorphisms, hence the fourth. The statement for pairs follows by the five lemma.
\end{proof}

\begin{corollary}\label{cor:euler}
The Betti numbers and the Euler characteristic of a finite simplicial complex depend only on the homotopy type of $|K|$. In particular all triangulations of a space have the same $V-E+F-\cdots$,
and for every convex polyhedron in $\R^3$, $V-E+F=2$.
\end{corollary}

\begin{proof}
By Theorem~\ref{thm:simplicialsingular}, $b_k$ is the rank of $H_k(|K|)$, and Proposition~\ref{prop:smithhomology} gives $\chi$ from the $b_k$. The boundary of a convex polyhedron, triangulated by cutting its faces into triangles
(which changes $E-F$ by nothing: each new diagonal adds one edge and one face), is homeomorphic to $S^2$ by radial projection from an interior point; so $\chi=b_0-b_1+b_2=1-0+1$.
\end{proof}

\begin{example}\label{ex:surfacesH}
\emph{Loops and $H_1$.} A loop $\gamma$ is a singular $1$-simplex with $\partial\gamma=0$. If $\gamma,\delta$ are loops at the same point, the singular $2$-simplex that runs $\gamma$ on $[e_0,e_1]$, $\delta$ on $[e_1,e_2]$
and $\gamma\cdot\delta$ on $[e_0,e_2]$ (compose $\gamma\cdot\delta$ with the affine map $\Delta^2\to[0,1]$ sending $e_0,e_1,e_2$ to $0,\frac12,1$) has boundary $\delta-\gamma\cdot\delta+\gamma$; and a homotopy of loops gives, by the
prism operator, a homology. So $[\gamma]\mapsto[\gamma]$ is a homomorphism $h\colon\pi_1(X,x_0)\to H_1(X)$. For the circle, Mayer--Vietoris with two arcs gives $H_1(S^1)\cong\tilde H_0(\text{two points})=\Z$,
and the loop $\omega_1$ maps to a generator (its two halves give $\delta[\omega_1]=\pm([p]-[q])$). So $h$ is the isomorphism $\pi_1(S^1)\cong H_1(S^1)$, and the degree of a map of the circle is its winding number.

\emph{Surfaces.} By induction with Mayer--Vietoris, $H_1$ of a wedge of $n$ circles is $\Z^n$, generated by the images of the $n$ loops, and $H_2=0$. Let $\Sigma_g$ be the polygon model of
Example~\ref{ex:surfaces}, $V$ a small open disc in the polygon and $U$ the complement of its centre. Then $U$ deformation retracts onto the wedge of $2g$ circles, $V$ is contractible, and $U\cap V$
deformation retracts onto a circle, whose generator maps in $H_1(U)$ to $h$ of the boundary word, $a_1+b_1-a_1-b_1+\dots=0$. Reduced Mayer--Vietoris reads
\[0\to H_2(\Sigma_g)\to H_1(S^1)=\Z\xrightarrow{0}\Z^{2g}\to H_1(\Sigma_g)\to\tilde H_0(S^1)=0,\]
so $H_2(\Sigma_g)=\Z$, $H_1(\Sigma_g)=\Z^{2g}$, $H_0=\Z$, and $\chi(\Sigma_g)=2-2g$. This is the value of the index sum of Chapter~4 for the pretzel with $g$ holes, which is homeomorphic to
$\Sigma_g$ (this follows from the classification of surfaces, a remark here).
\end{example}

\begin{remark}
Every compact surface admits a triangulation (Rad\'o, 1925), and every compact connected orientable surface is homeomorphic to exactly one $\Sigma_g$ (the classification, completed by Dehn and Heegaard,
1907, and Brahana, 1921). Together with Corollary~\ref{cor:euler} and Example~\ref{ex:surfacesH}, this identifies the index sum of Chapter~4 with $V-E+F$ for every compact orientable surface:
both are $2-2g$. A direct proof, building on each triangle a vector field with one zero at each vertex, edge and face, is in Hopf's 1926 paper.
\end{remark}

\begin{figure}[ht]
\centering
\fig{ch06-mv}
\caption{Mayer--Vietoris on the sphere: two open caps $U_+$ and $U_-$, each contractible, overlap in a band that retracts onto the equator. A cycle on $S^n$ is the sum of a chain in $U_+$
and a chain in $U_-$; its image under $\delta$ is their common boundary, a cycle in the band, one dimension lower. So every sphere carries exactly one hole, in its top dimension.}
\label{fig:mvsphere}
\end{figure}

\section{Invariance of domain and Borsuk--Ulam}

\begin{proposition}\label{prop:jordanbrouwer}
Let $n\ge1$. (a) If $h\colon\bar B^k\to S^n$ is an embedding (a continuous injection; a homeomorphism onto its image, $\bar B^k$ being compact), then $\tilde H_i(S^n\setminus h(\bar B^k))=0$ for all $i$.
(b) If $h\colon S^k\to S^n$ is an embedding with $k<n$, then $\tilde H_i(S^n\setminus h(S^k))$ is $\Z$ for $i=n-k-1$ and $0$ otherwise. In particular (Jordan--Brouwer) the complement of an embedded $S^{n-1}$
in $S^n$ has exactly two path components, and the complement of an embedded ball is path-connected.
\end{proposition}

\begin{proof}
(a) Replace $\bar B^k$ by the cube $I^k$, which is homeomorphic to it. All the complements below are nonempty, since the image of a cube is contractible and $S^n$ is not. Induction on $k$; for $k=0$, $S^n$ minus a point is $\R^n$. Suppose $\alpha\ne0$ in $\tilde H_i(S^n\setminus h(I^k))$. Let
$A=S^n\setminus h(I^{k-1}\times[0,\frac12])$ and $B=S^n\setminus h(I^{k-1}\times[\frac12,1])$, open sets with $A\cap B=S^n\setminus h(I^k)$ and $A\cup B=S^n\setminus h(I^{k-1}\times\{\frac12\})$, which has zero reduced homology by
induction. By Mayer--Vietoris, $\tilde H_i(A\cap B)\to\tilde H_i(A)\oplus\tilde H_i(B)$ is an isomorphism, so $\alpha$ stays nonzero in $A$ or in $B$. Repeating with the half-interval found, we get nested intervals
$I_m$ of length $2^{-m}$, shrinking to a point $p$, with $\alpha\ne0$ in $\tilde H_i(S^n\setminus h(I^{k-1}\times I_m))$ for every $m$. But by induction $\alpha$ is zero in $S^n\setminus h(I^{k-1}\times\{p\})$, the increasing union of
these open sets: a cycle representing $\alpha$ is the boundary of a chain there, and that chain, made of
finitely many simplices with compact images, lies in one $S^n\setminus h(I^{k-1}\times I_m)$. Contradiction.

(b) Induction on $k$. For $k=0$, $S^n$ minus two points is $\R^n$ minus a point, homotopy equivalent to $S^{n-1}$. For $k\ge1$, write $S^k=D_+\cup D_-$, two closed hemispheres meeting in $S^{k-1}$, and apply
Mayer--Vietoris to $A=S^n\setminus h(D_+)$ and $B=S^n\setminus h(D_-)$: their reduced homology vanishes by (a), $A\cap B=S^n\setminus h(S^k)$, $A\cup B=S^n\setminus h(S^{k-1})$, so
$\tilde H_i(S^n\setminus h(S^k))\cong\tilde H_{i+1}(S^n\setminus h(S^{k-1}))$. The last statements are the case $i=0$: an open subset of $S^n$ has path components that are open, and $\tilde H_0$ has rank one less than their number.
\end{proof}

\begin{theorem}[invariance of domain; Brouwer, 1912]\label{thm:invariancedomain}
If $U\subset\R^n$ is open and $h\colon U\to\R^n$ is continuous and injective, then $h(U)$ is open, and $h$ is a homeomorphism onto it.
\end{theorem}

\begin{proof}
Regard $\R^n$ as $S^n$ minus a point. Let $x\in U$ and $D\subset U$ a closed ball around $x$ with boundary sphere $S$. By Proposition~\ref{prop:jordanbrouwer}, $S^n\setminus h(S)$ has two path components, which are open.
It is the disjoint union of $S^n\setminus h(D)$, path-connected by (a), and $h(D\setminus S)$, path-connected as the image of an open ball. So these two sets are the two components, and $h(D\setminus S)$ is
open in $S^n$, contained in $h(U)$ and containing $h(x)$. So $h(U)$ is open, and $h$ maps open sets to open sets, so $h^{-1}$ is continuous.
\end{proof}

\begin{beauty}
Volume~6 promised this theorem, and it is worth pausing on how strange it is. Nothing in the hypothesis is about openness: $h$ is just a continuous injection. Yet the image of an open set must
be open, because a sphere drawn in the domain still separates space into an inside and an outside after $h$ has deformed it, and the inside is exactly where the ball went. The proof never
looks at $h$ closely; it only counts holes in complements. For a manifold, the theorem says that the dimension is a topological invariant: no injective continuous map can squash a
$3$-dimensional open set into a $2$-dimensional one, as Peano's space-filling curve might lead one to fear.
\end{beauty}

For the Borsuk--Ulam theorem we need coefficients $\Z/2$ and the projective spaces. Everything proved so far holds for any coefficients, with the same proofs.

\begin{lemma}\label{lem:rpn}
$H_k(\R P^n;\Z/2)=0$ for $k>n$.
\end{lemma}

\begin{proof}
Induction on $n$, with $\R P^0$ a point. Let $p=[0:\dots:0:1]$. On $U=\R P^n\setminus\{p\}$, $([x_0:\dots:x_n],t)\mapsto[x_0:\dots:x_{n-1}:(1-t)x_n]$ is a deformation retraction onto $\R P^{n-1}=\{x_n=0\}$ (it is
well defined because $x_0,\dots,x_{n-1}$ are not all zero). Let $V=\{[x]:|x_n|>\frac12|x|\}$, an open ball around $p$ in the chart $x_n=1$, and $U\cap V=V\setminus\{p\}\simeq S^{n-1}$. For $k>n$, the Mayer--Vietoris
sequence $H_k(U)\oplus H_k(V)\to H_k(\R P^n)\to H_{k-1}(U\cap V)$ has outer terms $0$ (by induction, contractibility of $V$, and $k-1>n-1$).
\end{proof}

\begin{theorem}[Borsuk--Ulam]\label{thm:borsukulam}
For $n\ge1$ there is no continuous map $g\colon S^n\to S^{n-1}$ with $g(-x)=-g(x)$. Equivalently, every continuous $f\colon S^n\to\R^n$ satisfies $f(x)=f(-x)$ for some $x$.
\end{theorem}

\begin{proof}
For $n=1$, $S^0$ is two points and $S^1$ is connected. Let $n\ge2$; all homology has coefficients $\Z/2$. Let $q\colon S^n\to\R P^n$ be the quotient map, a two-sheeted covering. Every singular simplex $\sigma$ in
$\R P^n$ has exactly two lifts to $S^n$ (Theorem~\ref{thm:liftcrit}: $\Delta^k$ is simply connected), $\tilde\sigma$ and $-\tilde\sigma$; let $\tau(\sigma)=\tilde\sigma+(-\tilde\sigma)$. Then $\tau$ is a chain map (the faces of the lifts are the lifts of the faces),
and
\[0\to C(\R P^n)\xrightarrow{\tau}C(S^n)\xrightarrow{q_\#}C(\R P^n)\to0\]
is exact: $q_\#$ is onto since every simplex lifts; $q_\#\tau=2=0$; and a chain $\sum a_{\tilde\sigma}\tilde\sigma$ is killed by $q_\#$ exactly when the two lifts of each $\sigma$ have equal coefficients, that is, when it is in the
image of $\tau$. The zigzag lemma gives the \emph{transfer sequence}
\[\cdots\to H_k(\R P^n)\xrightarrow{\tau_*}H_k(S^n)\xrightarrow{q_*}H_k(\R P^n)\xrightarrow{\delta}H_{k-1}(\R P^n)\to\cdots.\]
Now $H_k(S^n)=0$ for $0<k<n$, and $q_*$ is an isomorphism on $H_0=\Z/2$ (both spaces are path-connected). Reading the sequence: $\delta\colon H_k(\R P^n)\to H_{k-1}(\R P^n)$ is an isomorphism for
$1\le k\le n-1$ (at $k=1$, $\tau_*=0$ on $H_0$ because $q_*$ is injective there). At the top, $H_{n+1}(\R P^n)=0$ by Lemma~\ref{lem:rpn}, so $0\to H_n(\R P^n)\xrightarrow{\tau_*}H_n(S^n)=\Z/2\xrightarrow{q_*}H_n(\R P^n)\xrightarrow\delta H_{n-1}(\R P^n)\to0$
is exact. The sequence begins with $0$, so $\tau_*$ is injective. For $1\le k\le n-1$ the maps $\delta$ are isomorphisms and $H_0=\Z/2$, so $H_{n-1}(\R P^n)=\Z/2$. The next term after $H_{n-1}(\R P^n)$ is $H_{n-1}(S^n)=0$, and the sequence as written ends with $0$, so $\delta\colon H_n(\R P^n)\to H_{n-1}(\R P^n)$ is surjective. Let $A=H_n(\R P^n)$. Then $A$ injects into $\Z/2$, so $A$ is $0$ or $\Z/2$. If $A=0$, the map $\Z/2\xrightarrow{q_*}A$ cannot be exact at $\Z/2$: exactness would force $q_*$ to be injective, and $\Z/2$ does not inject into $0$. Thus $A=\Z/2$, the injection $\tau_*$ is an isomorphism, $q_*=0$, and $\delta$ is an isomorphism. Hence
$H_k(\R P^n)=\Z/2$ for $0\le k\le n$, with all these maps isomorphisms.

Suppose $g$ exists, and let $f\colon S^n\to S^n$ be $g$ followed by the inclusion of the equator. Then $f$ is odd, so it induces $\bar f\colon\R P^n\to\R P^n$, and $f_\#\tau=\tau\bar f_\#$ (the lifts of $\bar f\circ\sigma$ are
$f\circ\tilde\sigma$ and $f\circ(-\tilde\sigma)$). By naturality, $\bar f_*$ commutes with $\delta$ and $f_*\tau_*=\tau_*\bar f_*$. Now $\bar f_*$ is an isomorphism on $H_0$, hence on every $H_k(\R P^n)$, $k\le n$, by the isomorphisms
$\delta$; and then $f_*=\tau_*\bar f_*\tau_*^{-1}$ is an isomorphism on $H_n(S^n)$. But $f$ is not onto, so $f_*=0$ on $H_n(S^n)$. Contradiction.

Equivalence: if $f(x)\ne f(-x)$ for all $x$, then $g(x)=\frac{f(x)-f(-x)}{|f(x)-f(-x)|}$ is odd; conversely, an odd $g\colon S^n\to S^{n-1}\subset\R^n$ never satisfies $g(x)=g(-x)$.
\end{proof}

\begin{corollary}[the ham sandwich theorem]\label{cor:hamsandwich}
Let $A_1,\dots,A_n$ be bounded Lebesgue measurable subsets of $\R^n$. There is an affine hyperplane that cuts each of them into two halves of equal measure.
\end{corollary}

\begin{proof}
For $u=(u_0,u')\in S^n\subset\R\times\R^n$, let $H_u^+=\{x:\langle u',x\rangle>u_0\}$ and $f_i(u)=\lambda(A_i\cap H_u^+)$. For $u'\ne0$ the hyperplane $\langle u',x\rangle=u_0$ has measure zero, so $f_i$ is continuous by
dominated convergence (Volume~6); at $u=(\pm1,0)$, $H_u^+$ is empty or everything, and continuity there follows from the boundedness of $A_i$. By Theorem~\ref{thm:borsukulam} there is $u$ with $f_i(u)=f_i(-u)$ for all
$i$. Since $H_{-u}^+$ is the complement of $H_u^+$ up to a null set, $f_i(u)+f_i(-u)=\lambda(A_i)$, so $f_i(u)=\frac12\lambda(A_i)$. If $u=(\pm1,0)$ all the $A_i$ are null and any hyperplane works; otherwise
$\langle u',x\rangle=u_0$ is the hyperplane.
\end{proof}

Three measurable sets in space, the bread, the ham and the cheese, can be bisected by one plane cut. Hugo Steinhaus posed the problem in Lw\'ow, and
Stefan Banach solved it by reducing it to the Borsuk--Ulam theorem; a note in the Polish journal \emph{Mathesis Polska} recorded both in 1938, and Stone and Tukey proved the
$n$-dimensional version in 1942.

\section{Exercises}

\begin{exercise}[$\star$]\label{ex:kleinrp2}
\emph{The projective plane and the Klein bottle.} (a) The projective plane is a disc with its boundary glued by the word $aa$ (antipodal points). Show that $H_1(\R P^2)=\Z/2$ and $H_2(\R P^2)=0$, while
$H_2(\R P^2;\Z/2)=\Z/2$. (b) The Klein bottle is a square with sides glued by $aba^{-1}b$. Show that $H_1=\Z\oplus\Z/2$, $H_2=0$, and $\chi=0$.
\end{exercise}

\begin{exercise}[$\star\star$]\label{ex:hairydegree}
\emph{The hairy ball theorem by degree.} (a) Show that if $X$ is a tangent vector field on $S^n$ with no zero, then $(x,t)\mapsto\cos(\pi t)\,x+\sin(\pi t)\frac{X(x)}{|X(x)|}$ is a homotopy from the identity to the antipodal
map. (b) Deduce from Proposition~\ref{prop:reflection} that $n$ is odd, recovering Theorem~\ref{thm:hairyball} for continuous fields. (c) Show that a map $f\colon S^n\to S^n$ with no fixed point is homotopic to the
antipodal map, and that one with $f(x)\ne-x$ for all $x$ is homotopic to the identity. Deduce that if $n$ is even, every map $S^n\to S^n$ has a fixed point or sends some point to its antipode.
\end{exercise}

\begin{exercise}[$\star$]\label{ex:lhuilier}
\emph{Lhuilier's picture frame.} (a) A picture frame made of four square beams has an outer surface that can be triangulated as a torus. Check that $V-E+F=0$ on Figure~\ref{fig:torustriang}, and explain
with Corollary~\ref{cor:euler} why no triangulation of a torus gives $2$. (b) For a graph without loops or multiple edges (a $1$-dimensional simplicial complex) show that $b_0-b_1=V-E$, and compare with
Theorem~\ref{thm:wedge}.
\end{exercise}

\begin{exercise}[$\star\star$]\label{ex:coveringchi}
\emph{Euler characteristic of a covering.} Let $p\colon\tilde X\to|K|$ be a $k$-sheeted covering of a finite simplicial complex. (a) Show that each simplex of $K$ has exactly $k$ lifts, and that the lifts form a simplicial
complex triangulating $\tilde X$. (b) Deduce $\chi(\tilde X)=k\chi(K)$. (c) Deduce again that a subgroup of index $k$ in $F_n$ has rank $1+k(n-1)$, and that a $k$-sheeted covering of $\Sigma_g$ by a surface $\Sigma_h$ has
$2-2h=k(2-2g)$, as the remark of Chapter~5 announced.
\end{exercise}

\begin{exercise}[$\star\star$]\label{ex:lusternik}
\emph{Lusternik--Schnirelmann.} Show that if $S^n$ is covered by $n+1$ closed sets $F_1,\dots,F_{n+1}$, one of them contains a pair of antipodal points. (Use $f(x)=(d(x,F_1),\dots,d(x,F_n))$.)
\end{exercise}

\begin{exercise}[$\star$]\label{ex:pancakes}
\emph{Cutting points and pancakes.} (a) Two pancakes of any shapes lie on a plate. Show that one straight cut halves both. (b) Show that for any two finite sets of points in the plane there is a line such
that each open side contains at most half of each set. (Replace each point by a tiny disc and use (a); then shrink the discs.)
\end{exercise}

\section{Hints}

\begin{hint}{ex:kleinrp2}
Use the method of Example~\ref{ex:surfacesH}: the complement of a point retracts onto the $1$-skeleton, and the boundary circle maps to the boundary word in $H_1$ of the $1$-skeleton.
\end{hint}

\begin{hint}{ex:hairydegree}
For (c), push $f(x)$ towards $-x$ along the great circle from $f(x)$ that avoids $x$.
\end{hint}

\begin{hint}{ex:lhuilier}
For (a), count by rows. For (b), compute $H_0$ and $H_1$ of a graph from Theorem~\ref{thm:wedge}, or directly from Proposition~\ref{prop:smithhomology}.
\end{hint}

\begin{hint}{ex:coveringchi}
For (a), a simplex is simply connected, so it lifts; two lifts with a common vertex coincide by uniqueness of lifts.
\end{hint}

\begin{hint}{ex:lusternik}
Apply Borsuk--Ulam to $f$ and look at which $F_i$ contains $x$ and $-x$.
\end{hint}

\begin{hint}{ex:pancakes}
(a) is the ham sandwich theorem for $n=2$. For (b), a line meets each tiny disc in at most half its area unless it passes near the centre.
\end{hint}

\section{Solutions}

\begin{solution}{ex:kleinrp2}
\idea Glue a disc to a wedge of circles and its boundary becomes a relation; in $H_1$ the relation is the boundary word, written additively.
\writeup (a) The $1$-skeleton is a circle $a$, with $H_1=\Z$. The boundary of the disc runs around it twice, so in Mayer--Vietoris the map $H_1(S^1)=\Z\to H_1(a)=\Z$ is multiplication by $2$:
$0\to H_2(\R P^2)\to\Z\xrightarrow{\times2}\Z\to H_1(\R P^2)\to0$. Multiplication by $2$ is injective with cokernel $\Z/2$, so $H_2=0$ and $H_1=\Z/2$. With coefficients $\Z/2$ it is the zero map: $H_2=H_1=\Z/2$, as
Theorem~\ref{thm:borsukulam} found for all $\R P^n$. (b) The $1$-skeleton is a wedge of two circles $a,b$ (all four corners are identified), and the boundary word $aba^{-1}b$ gives $a+b-a+b=2b$. The map
$\Z\to\Z^2$, $1\mapsto(0,2)$, is injective with cokernel $\Z\oplus\Z/2$. So $H_2=0$, $H_1=\Z\oplus\Z/2$, and $\chi=1-1+0=0$.
\end{solution}

\begin{solution}{ex:hairydegree}
\idea A field with no zero gives every point a direction to walk in, and walking half a great circle in that direction takes every point to its antipode, continuously.
\writeup (a) Since $X(x)\perp x$, the point $\cos(\pi t)x+\sin(\pi t)\frac{X(x)}{|X(x)|}$ has norm $1$; it is $x$ at $t=0$ and $-x$ at $t=1$. (b) So $\deg(-\mathrm{id})=1$, while Proposition~\ref{prop:reflection} gives
$(-1)^{n+1}$; hence $n$ is odd. (c) If $f(x)\ne x$, the segment from $f(x)$ to $-x$ does not pass through $0$ (that would require $f(x)$ to be a positive multiple of $x$, so $f(x)=x$), and
$\frac{(1-t)f(x)-tx}{|(1-t)f(x)-tx|}$ is a homotopy from $f$ to $-\mathrm{id}$. Similarly, if $f(x)\ne-x$ then $\frac{(1-t)f(x)+tx}{|\cdots|}$ joins $f$ to $\mathrm{id}$. If $n$ is even, $\deg(-\mathrm{id})=-1\ne1=\deg\mathrm{id}$, so no map
can satisfy both conditions: either $f$ has a fixed point (else $\deg f=-1$) or it sends a point to its antipode (else $\deg f=1$).
\end{solution}

\begin{solution}{ex:lhuilier}
\idea Faces never matter, only holes: the torus has $b_0=1$, $b_1=2$, $b_2=1$.
\writeup (a) In Figure~\ref{fig:torustriang} each of the $9$ vertices of the torus is the lower left corner of one small square, which contributes three edges (bottom, left, diagonal) and two triangles, all
different after gluing: $V-E+F=9-27+18=0$. By Corollary~\ref{cor:euler}, $\chi$ is the same for all triangulations, and Example~\ref{ex:surfacesH} gives $1-2+1=0$: a solid whose surface is a torus, such
as Lhuilier's frame, can never satisfy Euler's formula. (b) Proposition~\ref{prop:smithhomology} with $c_k=0$ for $k\ge2$ gives $b_0-b_1=c_0-c_1=V-E$. For a connected graph $b_0=1$, so $b_1=E-V+1$, the rank of
the free group $\pi_1$ of Theorem~\ref{thm:wedge}.
\end{solution}

\begin{solution}{ex:coveringchi}
\idea Every simplex is simply connected, so it lifts to $k$ disjoint copies, and the counts $V$, $E$, $F$, \dots{} all multiply by $k$.
\writeup (a) A simplex $|\sigma|$ is simply connected and locally path-connected, so by Theorem~\ref{thm:liftcrit} its inclusion lifts with any prescribed value over one point: exactly $k$ lifts, one through each
point of the fibre over a vertex, and by Lemma~\ref{lem:uniquelift} two lifts of $\sigma$ that share a point coincide. The lifts of the faces of a lift are faces of that lift. Two lifted simplices with the same
vertices lie over the same simplex (whose vertex set they determine) and share a point, so they coincide. Order the lifted vertices compatibly with the order downstairs (break ties within a fibre
arbitrarily); the lifted simplices form a simplicial complex $\tilde K$, and the map $|\tilde K|\to\tilde X$ that sends each lifted simplex to its image is a continuous bijection from a compact space to a Hausdorff space,
a homeomorphism. (b) Each $c_j$ is multiplied by $k$. (c) For $R_n$ with a triangulation (subdivide each circle into three edges), $\chi=1-n$, and the covering graph has $\chi=1-\operatorname{rank}$; so
$1-\operatorname{rank}=k(1-n)$. For surfaces, $\chi(\Sigma_g)=2-2g$ by Example~\ref{ex:surfacesH}, and $\Sigma_g$ has a triangulation (subdivide the polygon model).
\end{solution}

\begin{solution}{ex:lusternik}
\idea Borsuk--Ulam gives a point that $f$ cannot tell from its antipode; then $x$ and $-x$ sit in the same first $n$ sets, or both outside all of them.
\writeup The map $f$ is continuous, so there is $x$ with $f(x)=f(-x)$. If $f_i(x)=0$ for some $i\le n$, then $d(x,F_i)=d(-x,F_i)=0$, and since $F_i$ is closed both $x$ and $-x$ lie in $F_i$. Otherwise neither $x$ nor $-x$
lies in $F_1\cup\dots\cup F_n$, so both lie in $F_{n+1}$.
\end{solution}

\begin{solution}{ex:pancakes}
\idea Every bisection problem with as many objects as dimensions is a ham sandwich.
\writeup (a) Corollary~\ref{cor:hamsandwich} with $n=2$. (b) Let $\varepsilon>0$ be less than half the smallest distance between two distinct points of $P\cup Q$, and replace each point by the open disc of radius
$\varepsilon$ around it; these discs are disjoint and have equal areas. By (a), some line $\ell_\varepsilon$ halves the total area of the discs of $P$ and that of the discs of $Q$. The discs of the points at
distance at least $\varepsilon$ from $\ell_\varepsilon$ lie entirely on one side, so each open side of $\ell_\varepsilon$ contains at most half of the points of $P$ at distance $\ge\varepsilon$ from it, and likewise for $Q$. Write
$\ell_\varepsilon=\{x:\langle u_\varepsilon,x\rangle=c_\varepsilon\}$ with $|u_\varepsilon|=1$; the line meets a fixed disc containing all the points, so the $c_\varepsilon$ are bounded, and for $\varepsilon=\frac1m$ a subsequence converges to
$(u,c)$, defining a line $\ell$. A point at positive distance from $\ell$ on one side lies, for $m$ large, on the same side of $\ell_{1/m}$ at distance more than $\frac1m$. So each open side of $\ell$ contains
at most half of $P$ and at most half of $Q$. Such discrete ham sandwich cuts are the basis of geometric divide-and-conquer algorithms; Lo, Matou\v sek and Steiger (1994) compute them in linear time in
the plane.
\end{solution}

\section*{Sources and further reading}
\addcontentsline{toc}{section}{Sources and further reading}

Euler's letter to Goldbach of 14 November 1750 is in \emph{Leonhard Euler und Christian Goldbach: Briefwechsel 1729--1764}, ed.\ A.~P.~Juškevič and E.~Winter (Akademie-Verlag, 1965); his proofs
are ``Elementa doctrinae solidorum'' and ``Demonstratio nonnullarum insignium proprietatum quibus solida hedris planis inclusa sunt praedita'', \emph{Novi Comm.\ Acad.\ Sci.\ Petrop.} 4 (1752/53, printed 1758),
109--140 and 140--160. S.~Lhuilier, ``M\'emoire sur la poly\'edrom\'etrie'', \emph{Ann.\ Math.\ Pures Appl.} 3 (1812--13), 169--189. I.~Lakatos, \emph{Proofs and Refutations} (Cambridge University Press, 1976);
D.~Richeson, \emph{Euler's Gem} (Princeton University Press, 2008), tells the whole story. On Emmy Noether: P.~Alexandroff and H.~Hopf, \emph{Topologie I} (Springer, 1935), preface; H.~Hopf, ``Eine
Verallgemeinerung der Euler--Poincar\'eschen Formel'', \emph{Nachr.\ Ges.\ Wiss.\ G\"ottingen} (1928), 127--136; C.~McLarty, ``Emmy Noether's set-theoretic topology'', in \emph{The Architecture of Modern
Mathematics} (Oxford University Press, 2006).

W.~Mayer, ``\"Uber abstrakte Topologie'', \emph{Monatsh.\ Math.\ Phys.} 36 (1929), 1--42 and 219--258; L.~Vietoris, ``\"Uber die Homologiegruppen der Vereinigung zweier Komplexe'', \emph{Monatsh.\ Math.\ Phys.}
37 (1930), 159--162. L.~E.~J.~Brouwer, ``Beweis der Invarianz des $n$-dimensionalen Gebiets'', \emph{Math.\ Ann.} 71 (1912), 305--313. K.~Borsuk, ``Drei S\"atze \"uber die $n$-dimensionale euklidische Sph\"are'',
\emph{Fund.\ Math.} 20 (1933), 177--190; A.~H.~Stone and J.~W.~Tukey, ``Generalized `sandwich' theorems'', \emph{Duke Math.\ J.} 9 (1942), 356--359; W.~A.~Beyer and A.~Zardecki, ``The early history of the ham sandwich theorem'', \emph{Amer.\ Math.\ Monthly} 111 (2004), 58--61, which translates the 1938 note. C.-Y.~Lo, J.~Matou\v sek and W.~Steiger, ``Algorithms for ham-sandwich cuts'', \emph{Discrete Comput.\ Geom.} 11 (1994), 433--452.

The proofs follow A.~Hatcher, \emph{Algebraic Topology} (Cambridge University Press, 2002), Sections~2.1--2.2 and~2.B (small chains, invariance of domain, and the transfer proof of Borsuk--Ulam),
with the comparison of simplicial and singular homology done by induction on simplices. J.~Matou\v sek, \emph{Using the Borsuk--Ulam Theorem} (Springer, 2003), is a delightful tour of the combinatorial
applications.
